% Source: https://artofproblemsolving.com/community/c3h1496341_polynomials
% Stable identity: aops-topic-1496341-q01
First, the rule holds for nonzero polynomials. If $p$ has leading term $a_mx^m$ and $q$ has leading term $b_nx^n$ with $a_m,b_n\ne0$, then the leading term of $pq$ is $a_mb_nx^{m+n}$, and $a_mb_n\ne0$. So $\deg(pq)=m+n$.

Now suppose $\deg0=d$ for some real number $d$, and take $p=0$ and $q=x$. Then $pq=0$, so the rule would give
\[
d=\deg(0\cdot x)=\deg0+\deg x=d+1 ,
\]
which is impossible. Hence no real number works. $\blacksquare$

With the convention $\deg0=-\infty$ and $-\infty+n=-\infty$, the rule holds in all cases. That is why the zero polynomial is usually given ``degree $-\infty$'' or no degree, while every nonzero constant has degree $0$.
